\documentclass[11pt]{article}
\usepackage[margin=1in]{geometry}
\usepackage{amsmath,amssymb,amsthm}
\usepackage{mathtools}
\usepackage{enumitem}
\usepackage{hyperref}
\usepackage{cleveref}
\usepackage{booktabs}
\usepackage{tabularx}
\usepackage{microtype}
\usepackage{cite}
\usepackage{xurl}
\usepackage{path}
\hypersetup{colorlinks=true,linkcolor=blue,citecolor=blue,urlcolor=blue}
\newtheorem{definition}{Definition}
\newtheorem{lemma}{Lemma}
\title{Request Fulfillment Certificates for Source-Only Releases\\\large Exact Completion Criteria for Public Request Contracts and Typed Omitted-Support Cuts}
\author{Anonymous}
\date{Unified archive note v0.12 (2026-03-21; archive v471)}
\begin{document}
\maketitle

\begin{abstract}
A terse source-only paper in this archive can now say exactly which maintained packet/menu cut stands behind its public ``supporting artifacts are available on request'' sentence, and it can also say what kind of omitted support that cut promises. One repeat ambiguity remains: when has such a request actually been satisfied? A responder may send only the public anchors, or only the validator bundle, or a broader directory missing one promised packet path. This note gives that completion criterion one home. It defines compact \emph{request fulfillment certificates}: family-scoped maintained objects saying which delivered paths, evidence-class slices, and validator companions constitute a complete response to one paper-facing source-only request promise.
\end{abstract}

\paragraph{Non-redundancy note.}
Synthesis~6 owns the archive-wide source-only artifact policy, Synthesis~52 owns exact answer disclosure packets, Synthesis~54 owns request-class review menus, Synthesis~56 owns exact paper-facing public request contracts, Synthesis~57 owns stable omitted-support family names, and Synthesis~17 owns the concrete worked instantiation.\cite{series:artifactpolicy,series:challengeanswerdisclosurepackets,series:challengeanswerreviewmenus,series:publicrequestcontracts,series:requestableevidenceclasses,series:workedexample}
This note owns only the completion criterion saying when a concrete response has fully satisfied one such source-only request promise. Imported evidence-class keys remain requestable-evidence vocabulary from Synthesis~57 rather than being redefined here, repeated family-completion verdict names may travel into later status envelopes, notices, routes, and spine mirrors only as imported family-scoped vocabulary, the within-release paper-level public token compressing those family verdicts belongs in Synthesis~59 rather than here, successor-release carryforward of that token belongs in Synthesis~60 rather than here, the exact outward-facing notice sentence readers should see belongs in Synthesis~61 rather than here, and the canonical sentence family plus slot realization rule behind that wrapper belong in Synthesis~62 rather than here, the request-side terminal stop belongs in Synthesis~73 rather than here, the paired answer/request citation discipline belongs in Synthesis~74 rather than here, and the stop-versus-widen/reopen rule belongs in Synthesis~75 rather than here.

\paragraph{Scope.}
This is not a new disclosure packet, review menu, validator, or release stage.
It is one thin completion layer saying when a paper-facing request promise has been fully discharged for one public-anchor family.

\paragraph{Exports.}
This note exports (i) request fulfillment certificates, (ii) a family-completion rule, (iii) a validator-companion completion rule, (iv) an imported-evidence-class-vocabulary rule, (v) a shared-class discharge rule, (vi) a shared-completion-vocabulary rule, and (vii) a contract-satisfaction rule saying that later papers may cite one compact maintained object when they need an auditable answer to ``has the source-only request actually been fulfilled?'' Later terse papers should cite Synthesis~58 only when the live issue is the exact family-scoped completion basis itself: whether one promised family has been fully discharged, which delivered support paths and validator companions made that true, or which family-scoped completion verdict is the semantic owner of one later imported status token or route-side mirror. They should stop at Synthesis~56 when the live issue is the paper-facing request promise itself, stop at Synthesis~57 when the live issue is the omitted-support family typing or shared evidence-class vocabulary, widen to Synthesis~59--60 when the live issue has become the within-release paper-level status token or its cross-release carryforward, stop at Synthesis~61--63 when the completion basis is already fixed and the remaining issue is the exact outward-facing sentence, canonical notice family, or token-keyed family selector shown to readers, widen through Synthesis~64--72 only when that visible-wrapper layer is already fixed and the live issue has moved into the follow-up handoff, predecessor packet cut, successor reuse verdict, packet-local refresh delta, visible refresh sentence, canonical refresh family, family-selection judgment, successor-facing handoff, or exact successor-facing packet cut, widen to Synthesis~73 when the live issue is whether that request-side branch already closes, widen to Synthesis~74--75 when the remaining issue is paired-terminal reuse or fail-closed cutover, and stop at Synthesis~17 when the concrete maintained worked route is the live issue instead.\cite{series:publicrequestcontracts,series:requestableevidenceclasses,series:publicrequeststatusenvelopes,series:publicrequestcarryforward,series:publicrequestnotices,series:publicrequestnoticenormalforms,series:publicrequestnoticeselections,series:publicrequestresponsemenus,series:publicrequestresponsepackets,series:publicrequestresponsepacketcarryforward,series:publicrequestresponsepacketdeltas,series:publicrequestrefreshnotices,series:publicrequestrefreshnormalforms,series:publicrequestrefreshselections,series:publicrequestsuccessormenus,series:publicrequestsuccessorpackets,series:requestsideterminalcitationnormalform,series:pairedterminalcitationdiscipline,series:pairedterminalcutoverrules,series:workedexample}

\section{Why request fulfillment certificates deserve their own note}
Public request contracts say what a short source-only paper is promising.
Requestable evidence classes say what kind of omitted-support family stands behind that promise.
Disclosure packets and review menus say which maintained packet/menu objects govern the handoff.
A maintainer or referee still needs one more stable question answered without prose drift: when is the response complete?
Without one maintained completion object, one responder may send the packet slice but omit validator companions, another may send the validator bundle but omit the family-specific support paths, and a third may send a larger archive cut without proving that it actually covers the promised maintained slice.
That is the same old ambiguity, just shifted from promise formation to promise satisfaction.

\begin{definition}[Request fulfillment certificate]
Fix one public request contract key $r$ and one covered public-anchor family key $f$.
A \emph{request fulfillment certificate} is a tuple
\[
H(r,f)=(r,f,\Pi,E,D,V,\nu,\mu)
\]
containing one imported disclosure-packet key set $\Pi$, one imported evidence-class key set $E$, one delivered support-path set $D$, one validator-companion path set $V$, one completion verdict $\nu\in\{\texttt{complete},\texttt{incomplete}\}$, and one verdict-posture tag $\mu$.
The certificate does not create a new packet or contract.
It only states the exact completion cut for one already-promised source-only request family.
\end{definition}

\begin{definition}[Family completion]
A request fulfillment certificate satisfies the \emph{family-completion rule} if its delivered support-path set covers every request-packet path named by the imported disclosure packets for that family.
It must also cover every canonical support path named by the imported evidence classes for that same family.
So the certificate does not accept a vague broader-answer claim in place of the maintained cut.
\end{definition}

\begin{definition}[Validator-companion completion rule]
A request fulfillment certificate satisfies the \emph{validator-companion completion rule} if its validator-companion path set contains the validation companions required by the imported public request contract and keeps those companions distinct from the family-specific delivered support paths.
So completion preserves both the evidence slice and the recheck hook.
\end{definition}

\begin{definition}[Imported-evidence-class-vocabulary rule]
A request fulfillment certificate satisfies the \emph{imported-evidence-class-vocabulary rule} if every evidence-class key it repeats remains imported requestable-evidence-class vocabulary: the certificate may say which shared and family-specific classes it is discharging, but it does not thereby redefine the canonical support-path meaning of those class keys.
So completion may point back to stable class names without becoming the owner of the omitted-support family definitions those names denote.
\end{definition}
\begin{definition}[Shared-class discharge rule]
A request fulfillment certificate satisfies the \emph{shared-class discharge rule} if every shared requestable evidence class it imports is discharged only \emph{as imported into that certificate's family}: a class reused across multiple families may therefore appear in multiple family-scoped certificates without forcing a separate family-agnostic completion verdict.
So shared omitted-support vocabulary does not collapse the completion boundary that still lives at the family level.
\end{definition}

\begin{definition}[Shared-completion-vocabulary rule]
A request fulfillment certificate satisfies the \emph{shared-completion-vocabulary rule} if its completion verdict name may recur in later status envelopes, visible notices, routes, or spine mirrors only as imported family-scoped vocabulary: the same verdict label may be reused across many paper-level or wrapper objects, but the truth of that verdict remains owned by the concrete family certificate that emits it.
So later compression layers may import the verdict name without becoming the completion owner.
\end{definition}

\begin{definition}[Contract-satisfaction rule]
A public request contract satisfies the \emph{contract-satisfaction rule} if every public-anchor family it covers has at least one companion request fulfillment certificate satisfying the family-completion rule, the validator-companion completion rule, the imported-evidence-class-vocabulary rule, the shared-class discharge rule, and the shared-completion-vocabulary rule, and if every such certificate imports only disclosure-packet and evidence-class keys already covered by that contract.
So the contract keeps ownership of the promise while the certificate owns only the completion verdict for one family.
\end{definition}

\begin{lemma}[Request fulfillment certificates stabilize source-only completion without widening promise ownership]
If a public request contract is accompanied by request fulfillment certificates satisfying the family-completion, validator-companion completion, imported-evidence-class-vocabulary, shared-class discharge, shared-completion-vocabulary, and contract-satisfaction rules, then a later terse paper may cite those certificates to show that its source-only request promise has been fully discharged without changing the ownership of packet paths, review-menu choices, evidence-class naming, or validator semantics.
\end{lemma}

\begin{proof}
The family-completion rule prevents the completion object from shrinking the promised support cut below the imported disclosure packets and evidence classes.
The validator-companion completion rule prevents the completion object from treating bare evidence paths as enough when the promise also required a maintained recheck hook. The imported-evidence-class-vocabulary rule closes the upstream seam: even when the same evidence-class keys recur inside completion certificates and later route/spine mirrors, the owner of the class meaning remains the evidence-class ledger rather than the completion layer.
The shared-class discharge rule closes the remaining seam for shared evidence classes: even when one evidence class is reused across multiple families, completion is still recorded through the family-scoped certificate that imports it rather than through a separate supra-family verdict. The shared-completion-vocabulary rule closes the next seam: even when the same completion verdict name is mirrored into later status envelopes, notices, routes, and spine surfaces, the semantic owner of that verdict remains the concrete family-scoped certificate. The contract-satisfaction rule prevents the completion object from inventing new families or new packet ownership outside the imported paper-facing contract.
So the certificate changes only status: it records whether a promised cut has been fully delivered, but it does not replace the packet, menu, contract, evidence-class, or validator owners.\qedhere
\end{proof}

\paragraph{A forced family-completion matrix.}
\begin{table}[t]
\centering
\footnotesize
\begin{tabularx}{\linewidth}{@{}p{0.18\linewidth}p{0.18\linewidth}p{0.16\linewidth}p{0.16\linewidth}X@{}}
\toprule
Case & Imported packet/class cut covered? & Validator companions present? & Shared class discharged in this family certificate? & Forced family verdict / honest stop \\
\midrule
Dropped family path & No & either & either & \texttt{incomplete}. Reopen Synthesis~58 immediately because the certificate no longer covers the promised family cut. \\
Validator gap & Yes & No & Yes & \texttt{incomplete}. The evidence slice may be present, but the contract's recheck hook is missing, so completion fails at the validator-companion rule. \\
Borrowed sibling discharge & Yes & Yes & No & \texttt{incomplete}. A shared class satisfied only through another family's certificate does not discharge this family; the completion boundary stays family-scoped. \\
Exact family discharge & Yes & Yes & Yes & \texttt{complete}. Later terse papers may stop at the family certificate unless the live issue widens to paper-level status, visible notice, or request-side terminal reuse. \\
\bottomrule
\end{tabularx}
\caption{A forced family-completion matrix. Completion is not a prose impression: for one promised family it is forced by three local checks --- the imported packet/class cut is covered, the validator companions required by the contract are present, and any shared class is discharged inside this family's own certificate rather than borrowed from a sibling family.}
\label{tab:forced-family-completion-matrix}
\end{table}

This matrix is intentionally non-decorative: it makes explicit that there is no supra-family shortcut from ``some sibling family is complete'' to ``this family is complete,'' and no packet-only shortcut from ``the support files arrived'' to ``the request promise is fully discharged.'' Completion remains a forced family-scoped verdict owned by the concrete certificate that names the delivered paths and validator companions.\cite{series:workedexample}

\paragraph{One tiny force drill.}
In the worked note, both public-anchor families import the same shared class \texttt{bundle\_identity\_and\_validator\_companions}, but they still earn \texttt{completion\_verdict = complete} separately. The compact-diff certificate is not allowed to point at the public-status certificate and say ``the validator class was already discharged over there.'' It must carry that shared class in its own imported class set, keep the shared validator companions in its own \texttt{validation\_paths}, and keep its own family-local compare/replay support cut. So the family-scoped completion question stays auditable: if the compact-diff certificate lost \texttt{example\_validation\_report.json}, or if it tried to rely on the sibling family's certificate instead of importing the shared class itself, the honest verdict would become \texttt{incomplete} even though another family in the same paper remained complete.\cite{series:workedexample}

\paragraph{A minimal public request-fulfillment card.}
\begin{table}[t]
\centering
\small
\begin{tabularx}{0.97\linewidth}{@{}p{0.31\linewidth}X@{}}
\toprule
\textbf{Public row} & \textbf{Pinned value}\\
\midrule
Contract anchor & \texttt{public\_request\_contract\_key = worked\_example\_receipt\_interlock\_public\_request\_contract} from \nolinkurl{example_public_request_contracts.json} \\
Family cut & \texttt{family\_key = compact\_diff\_summary\_family} with certificate key \texttt{compact\_diff\_summary\_request\_fulfillment\_certificate} \\
Imported disclosure / class floor & The certificate covers \texttt{release\_note\_compact\_diff\_summary\_answer\_disclosure\_packet} and \texttt{auditor\_compact\_diff\_summary\_answer\_disclosure\_packet}, imports the shared class \texttt{bundle\_identity\_and\_validator\_companions}, and imports the family-local class \texttt{compact\_diff\_compare\_and\_replay\_support} \\
Delivered public floor & \texttt{delivered\_public\_paths = [example\_compare\_report.json, example\_successor\_clause\_pack.json]} \\
Delivered support floor & \texttt{delivered\_support\_paths} extends that public floor with the family-local compare/replay slice, including \texttt{example\_successor\_sentence\_locks.json}, \texttt{example\_verifier\_report.json}, and \texttt{example\_successor\_continuity\_verdict.json} \\
Validation companions & \texttt{validation\_paths = [support\_manifest.json, example\_artifact\_inventory.json, example\_validation\_report.json, ../tools/validate\_example.py]} \\
Completion verdict & \texttt{completion\_verdict = complete} with \texttt{completion\_verdict\_mode = shared\_vocabulary\_family\_scoped\_verdict} \\
Escalation pointer & Widen to Synthesis~56 for the paper-facing promise, Synthesis~57 for evidence-class semantics, Synthesis~59--63 for paper-level status / notices, and Synthesis~73--75 once the live issue is request-side terminal reuse or cutover \\
\bottomrule
\end{tabularx}
\caption{A minimal public request-fulfillment card. This is the smallest public answer later terse papers can quote directly when the live issue is whether one promised source-only family was fully discharged, by which concrete delivered paths and validator companions, without reopening the wider paper-level status or request-side routing chain.}
\label{tab:minimal-request-fulfillment-card}
\end{table}

This card is intentionally non-decorative: it pins one contract key, one family key, one imported disclosure/class floor, one delivered public/support cut, one validator-companion cut, and one family-scoped completion verdict without silently re-owning the paper-level status token, the visible notice sentence, or the request-side terminal stop.\cite{series:workedexample}

\paragraph{One tiny completion drill.}
In the worked compact-diff family, the public request contract promises both the terse public compare anchors and the maintained compare/replay support slice. The fulfillment certificate is therefore honest only because its delivered paths include the public anchors \texttt{example\_compare\_report.json} and \texttt{example\_successor\_clause\_pack.json}, its support cut includes the family-local compare/replay artifacts, and its validation cut still includes the shared bundle-identity companions. That is why the family may publish \texttt{completion\_verdict = complete}. If the same certificate kept the support slice but dropped \texttt{example\_validation\_report.json} or \nolinkurl{../tools/validate_example.py}, the certificate would no longer satisfy the validator-companion rule and later terse papers would have to reopen Synthesis~58 rather than quote this compact card as complete.\cite{series:workedexample}

\section{Worked example pointer}
The maintained worked bundle now ships \nolinkurl{example_request_fulfillment_certificates.json}.\cite{series:workedexample}
It records two family-scoped fulfillment certificates for the worked note: one for the public-status lineage/closure family and one for the compact-diff compare/replay family.
Each certificate imports the paper-facing public request contract, splits its evidence classes into the shared bundle-identity/validator companion class and the family-local support class, carries the delivered support-path cut, names the validator companions, and records one completion verdict. The worked bundle now also mirrors the imported evidence-class vocabulary posture directly, so later route/spine surfaces can repeat those class keys through \nolinkurl{source_only_evidence_class_mode_map} without pretending that completion owns class meaning.
So a later terse paper can say not only what its source-only sentence promises, but also what exact maintained cut counts as a complete response once a referee asks for the omitted support, even when one evidence class is shared across families. The worked example now also re-exports those same certificate keys through its answer routes and packet-local family-basis maps, so downstream terse notes can recover completion without detouring through the whole source-only tail or reopening the exact packet ledgers just to learn which completion families a bounded handoff slice is carrying. Under the paired terminal citation discipline and paired cutover rule, later terse papers should usually stop at Synthesis~73--75 for the route-side stop/widen judgment and reopen these certificates only when the exact family-scoped discharge question is itself disputed. The paper-level public token compressing those family verdicts then belongs in Synthesis~59, the successor-release carryforward answer belongs in Synthesis~60, and the exact outward-facing sentence readers should see belongs in Synthesis~61.\cite{series:requestsideterminalcitationnormalform,series:pairedterminalcitationdiscipline,series:pairedterminalcutoverrules}

\begin{thebibliography}{99}\small

\bibitem{series:artifactpolicy}
Anonymous.
\newblock Artifact and Disclosure Policy for the One-True-Archive (Source-Only Public Release).
\newblock Series synthesis note (Synthesis~6), 2026. (This archive.)

\bibitem{series:challengeanswerdisclosurepackets}
Anonymous.
\newblock Challenge Answer Disclosure Packets for Successor Maintenance: Public Pointers and Requestable Audit Slices under Space Pressure.
\newblock Series synthesis note (Synthesis~52), 2026. (This archive.)

\bibitem{series:challengeanswerreviewmenus}
Anonymous.
\newblock Challenge Answer Review Menus for Successor Maintenance: Smallest-Sufficient Referee Handoffs from Dockets, Profiles, Packets, and Ladders.
\newblock Series synthesis note (Synthesis~54), 2026. (This archive.)

\bibitem{series:publicrequestcontracts}
Anonymous.
\newblock Public Request Contracts for Source-Only Releases: Exact Paper-Level Handoffs from Boilerplate Policy to Maintained Packet Cuts.
\newblock Series synthesis note (Synthesis~56), 2026. (This archive.)

\bibitem{series:requestableevidenceclasses}
Anonymous.
\newblock Requestable Evidence Classes for Source-Only Releases: Stable Names for Omitted-Support Families behind Packet Cuts and Public Request Contracts.
\newblock Series synthesis note (Synthesis~57), 2026. (This archive.)


\bibitem{series:publicrequeststatusenvelopes}
Anonymous.
\newblock Public Request Status Envelopes for Source-Only Releases: Paper-Level Tokens Compressing Family Completion without Re-owning the Promise.
\newblock Series synthesis note (Synthesis~59), 2026. (This archive.)

\bibitem{series:publicrequestcarryforward}
Anonymous.
\newblock Public Request Carryforward for Source-Only Successor Releases: When Paper-Level Request Status Persists, Advances, or Reopens across Release Evolution.
\newblock Series synthesis note (Synthesis~60), 2026. (This archive.)

\bibitem{series:publicrequestnotices}
Anonymous.
\newblock Public Request Notices for Source-Only Releases: Outward-Facing Sentences for Within-Release Status and Cross-Release Carryforward.
\newblock Series synthesis note (Synthesis~61), 2026. (This archive.)

\bibitem{series:publicrequestnoticenormalforms}
Anonymous.
\newblock Public Request Notice Normal Forms for Source-Only Releases: Canonical Sentence Families for Within-Release Status and Cross-Release Carryforward Notices.
\newblock Series synthesis note (Synthesis~62), 2026. (This archive.)

\bibitem{series:publicrequestnoticeselections}
Anonymous.
\newblock Public Request Notice Selections for Source-Only Releases: Selecting One Canonical Public Wrapper from Stable Status and Carryforward Families.
\newblock Series synthesis note (Synthesis~63), 2026. (This archive.)

\bibitem{series:publicrequestresponsemenus}
Anonymous.
\newblock Public Request Response Menus for Source-Only Releases: Smallest-Sufficient Maintained Follow-Ups for Public Notice Questions.
\newblock Series synthesis note (Synthesis~64), 2026. (This archive.)

\bibitem{series:publicrequestresponsepackets}
Anonymous.
\newblock Public Request Response Packets for Source-Only Releases: Exact Bounded Handoffs from Public Notice Questions to Maintained Omitted-Support Slices.
\newblock Series synthesis note (Synthesis~65), 2026. (This archive.)

\bibitem{series:publicrequestresponsepacketcarryforward}
Anonymous.
\newblock Public Request Response Packet Carryforward for Source-Only Successor Releases: Reuse, Refresh, and Reopen Rules for Bounded Omitted-Support Handoffs.
\newblock Series synthesis note (Synthesis~66), 2026. (This archive.)

\bibitem{series:publicrequestresponsepacketdeltas}
Anonymous.
\newblock Public Request Response Packet Delta Ledgers for Source-Only Successor Releases: Exact Refresh Deltas for Reused Omitted-Support Families.
\newblock Series synthesis note (Synthesis~67), 2026. (This archive.)

\bibitem{series:publicrequestrefreshnotices}
Anonymous.
\newblock Public Request Response Packet Refresh Notices for Source-Only Successor Releases: Visible Reissue Sentences after Packet Refresh.
\newblock Series synthesis note (Synthesis~68), 2026. (This archive.)

\bibitem{series:publicrequestrefreshnormalforms}
Anonymous.
\newblock Public Request Response Packet Refresh Notice Normal Forms for Source-Only Successor Releases: Canonical Visible Wrapper Families after Packet Refresh.
\newblock Series synthesis note (Synthesis~69), 2026. (This archive.)

\bibitem{series:publicrequestrefreshselections}
Anonymous.
\newblock Public Request Response Packet Refresh Notice Selections for Source-Only Successor Releases: Choosing One Canonical Visible Refresh Family from Stable Packet Deltas.
\newblock Series synthesis note (Synthesis~70), 2026. (This archive.)

\bibitem{series:publicrequestsuccessormenus}
Anonymous.
\newblock Public Request Response Packet Refresh Response Menus for Source-Only Successor Releases: Smallest-Sufficient Successor Handoffs after Visible Refresh.
\newblock Series synthesis note (Synthesis~71), 2026. (This archive.)

\bibitem{series:publicrequestsuccessorpackets}
Anonymous.
\newblock Public Request Response Packet Refresh Response Packets for Source-Only Successor Releases: Exact Successor-Facing Packet Cuts after Visible Refresh.
\newblock Series synthesis note (Synthesis~72), 2026. (This archive.)

\bibitem{series:requestsideterminalcitationnormalform}
Anonymous.
\newblock Public Request Response Packet Refresh Response Packet Closure Verdicts for Source-Only Successor Releases: Terminal Closure Stops for the Visible Refresh Branch.
\newblock Series synthesis note (Synthesis~73), 2026. (This archive.)

\bibitem{series:pairedterminalcitationdiscipline}
Anonymous.
\newblock Paired Terminal Citation Discipline for Stable Review Spines: One Answer-Side Stop and One Request-Side Capstone for Terse Papers.
\newblock Series synthesis note (Synthesis~74), 2026. (This archive.)

\bibitem{series:pairedterminalcutoverrules}
Anonymous.
\newblock Paired Terminal Cutover Rules for Stable Review Spines: When to Stop at the Checked Answer, When to Widen to Source-Only, and When to Fail Closed to the Reopening Owner.
\newblock Series synthesis note (Synthesis~75), 2026. (This archive.)

\bibitem{series:workedexample}
Anonymous.
\newblock Worked Example: Receipt Interlock for an Anonymous-DHT Reader-Privacy Claim.
\newblock Series synthesis note (Synthesis~17), 2026. (This archive.)

\end{thebibliography}
\end{document}
